\chapter{Puissances}\label{ch:g8:powers}

Plier une feuille de papier en deux $10$ fois (si l'on peut~!)~: son
épaisseur double à chaque fois, et $10$ doublements la multiplient par
$2^{10} = 1024$. Les puissances sont l'abréviation de la
multiplication répétée~; ce chapitre fixe la notation et ses règles,
avec un rôle particulier pour les puissances de $10$.

\section{Définition}

\begin{definition}[Puissance]\label{def:g8:powers:def}
Pour un nombre $a$ et un entier $n \geq 1$~:
\[
a^n = \underbrace{a \times a \times \dots \times a}_{n
\text{ facteurs}},
\]
lu «~$a$ à la (puissance) $n$~»~; $a$ est la \emph{base}, $n$
l'\emph{exposant}\index{exposant}. Noms particuliers~: $a^2$ se lit
«~$a$ au carré~», $a^3$ «~$a$ au cube~». Par convention~:
\[
a^1 = a,
\qquad
a^0 = 1 \ (a \neq 0),
\qquad
a^{-n} = \frac{1}{a^n}.
\]
\end{definition}

\begin{example}\label{ex:g8:powers:def}
$3^4 = 3 \times 3 \times 3 \times 3 = 81$~; \quad
$10^3 = 1000$~; \quad $5^{-2} = \frac{1}{25} = 0.04$~; \quad
$(-2)^3 = -8$ et $(-2)^4 = +16$
(\cref{ex:g8:negprod:powers}). Attention~: $a^n$ n'est \emph{pas}
$a \times n$~: \ $2^5 = 32$, et non $10$.
\end{example}

\begin{omfigure}
\begin{tikzpicture}[scale=0.9]
  \foreach \n/\x in {0/0, 1/1.4, 2/2.8, 3/4.2, 4/5.6} {
    \pgfmathsetmacro{\h}{0.28*pow(2,\n)}
    \draw[omDef, fill=omDef!20] (\x,0) rectangle (\x+0.9,\h);
    \node[below, font=\small] at (\x+0.45,0) {$2^\n$};
    \pgfmathtruncatemacro{\v}{pow(2,\n)}
    \node[above, font=\small] at (\x+0.45,\h) {$\v$};
  }
\end{tikzpicture}

{\small Puissances de $2$~: chaque barre est le double de la
précédente. La croissance par multiplication répétée s'emballe bien
plus vite que la croissance par addition répétée.}
\end{omfigure}

\section{Les règles des exposants}

\begin{theorem}[Règles des exposants]\label{thm:g8:powers:rules}
Pour une base non nulle $a$ et des \omterm{def:g8:powers:def}{exposants} entiers $m$, $n$~:
\[
a^m \times a^n = a^{m+n},
\qquad
\frac{a^m}{a^n} = a^{m-n},
\qquad
\left(a^m\right)^n = a^{m \times n},
\qquad
(ab)^n = a^n b^n .
\]
\end{theorem}

\begin{proof}[Preuve en comptant les facteurs]
$a^m \times a^n$ aligne $m$ facteurs $a$ suivis de $n$ autres~:
$m + n$ au total. $\left(a^m\right)^n$ répète un bloc de $m$ facteurs
$n$ fois~: $mn$ facteurs. $(ab)^n$ contient $n$ lettres $a$ et $n$
lettres $b$, que l'on peut regrouper. La règle du \omterm{def:g3:division:remainder}{quotient} vient
d'annuler $n$ des $m$ facteurs~; avec les conventions $a^0 = 1$ et
$a^{-n} = \frac{1}{a^n}$, elle reste vraie même lorsque $n \geq m$.
\end{proof}

\begin{example}\label{ex:g8:powers:rules}
\[
2^3 \times 2^5 = 2^8 = 256,
\qquad
\frac{7^6}{7^4} = 7^2 = 49,
\qquad
\left(5^2\right)^3 = 5^6,
\qquad
2^4 \times 5^4 = 10^4 .
\]
Un piège~: les règles s'appliquent à une \emph{base commune} (ou à un
\omterm{def:g8:powers:def}{exposant} commun, pour la dernière). Aucune règle ne simplifie
$2^3 \times 5^2$ --- il faut juste calculer~: $8 \times 25 = 200$.
\end{example}

\section{Puissances de dix}

\begin{proposition}[Puissances de dix]\label{prop:g8:powers:ten}
Pour $n \geq 1$~: $10^n = 1\underbrace{0\dots0}_{n}$, et
$10^{-n} = 0.\underbrace{0\dots0}_{n-1}1$. Les règles des \omterm{def:g8:powers:def}{exposants}
se lisent~: multiplier des puissances de dix, c'est additionner les
\omterm{def:g8:powers:def}{exposants}.
\end{proposition}

\begin{example}\label{ex:g8:powers:ten}
$10^4 \times 10^3 = 10^7$~; \quad
$\dfrac{10^2}{10^5} = 10^{-3} = 0.001$. Les grandes et petites
quantités deviennent lisibles~:
\[
\text{un milliard} = 10^9,
\qquad
\text{un millionième} = 10^{-6}.
\]
Combiné avec les décimaux~: $3.2 \times 10^5 = 320\,000$ et
$4.7 \times 10^{-3} = 0.0047$ (décaler la \omterm{def:g4:decimals:point}{virgule},
\cref{prop:g6:decimals:shift}). L'usage systématique de cette
écriture --- la \emph{notation scientifique} --- est développé au
\cref{ch:g9:fractions}.
\end{example}

\begin{example}[Ordres de grandeur]\label{ex:g8:powers:magnitude}
La lumière parcourt environ $3 \times 10^8$ m/s~; une année compte
environ $3.2 \times 10^7$ secondes. Une année-lumière mesure donc
environ
\[
3 \times 3.2 \times 10^{8+7} \approx 10 \times 10^{15} = 10^{16}
\text{ m}
\]
--- dix millions de milliards de mètres. Les puissances de dix font
tenir les calculs astronomiques sur une ligne.
\end{example}

\begin{method}[Simplifier une expression avec des puissances]\label{met:g8:powers:simplify}
\begin{enumerate}
  \item Grouper les facteurs base par base~;
  \item appliquer les règles d'\omterm{def:g8:powers:def}{exposants} au sein de chaque base~;
  \item calculer les petites puissances restantes, ou laisser la
        réponse sous forme de puissance si elle est grande.
\end{enumerate}
\end{method}

\begin{example}\label{ex:g8:powers:simplify}
\[
\frac{3^5 \times 3^{-2} \times 4^2}{3^2}
= \frac{3^{5 + (-2)}}{3^2} \times 16
= 3^{3 - 2} \times 16
= 3 \times 16 = 48 .
\]
\end{example}

\section{Exercices}

\begin{exercise}[$\star$]\label{exo:g8:powers:1}
Calculer~:
\[
2^6, \qquad 3^3, \qquad 10^5, \qquad 1^{100}, \qquad 0.1^2, \qquad
6^0 .
\]
\end{exercise}

\begin{exercise}[$\star$]\label{exo:g8:powers:2}
Calculer~:
\[
(-3)^2, \qquad -3^2, \qquad (-1)^{15}, \qquad (-5)^3, \qquad
\left(\tfrac{2}{3}\right)^2 .
\]
\end{exercise}

\begin{exercise}[$\star$]\label{exo:g8:powers:3}
Écrire comme une seule puissance~:
\[
7^4 \times 7^5, \qquad
\frac{2^9}{2^3}, \qquad
\left(10^3\right)^4, \qquad
5^6 \times 5^{-2}, \qquad
3^4 \times 7^4 .
\]
\end{exercise}

\begin{exercise}[$\star$]\label{exo:g8:powers:4}
Écrire comme un \omterm{def:g6:decimals:places}{nombre décimal}~: $10^{-2}$~; \quad $4 \times 10^3$~;
\quad $2.5 \times 10^{-4}$~; \quad $10^0$.
\end{exercise}

\begin{exercise}[$\star$]\label{exo:g8:powers:5}
Écrire avec une puissance de dix~: cent mille~; un dixième~; dix
milliards~; $0.000\,001$.
\end{exercise}

\begin{exercise}[$\star$]\label{exo:g8:powers:6}
Simplifier, puis calculer~:
\[
\frac{10^7 \times 10^{-3}}{10^2},
\qquad
\frac{2^5 \times 2^4}{2^6},
\qquad
\left(2^2\right)^3 \times 2^{-4} .
\]
\end{exercise}

\begin{exercise}[$\star$]\label{exo:g8:powers:7}
Vrai ou faux~? Corriger les faux.
\[
2^3 \times 2^4 = 2^{12}; \qquad
5^2 + 5^3 = 5^5; \qquad
(3^2)^4 = 3^8; \qquad
10^3 \times 10^3 = 100^3 .
\]
\end{exercise}

\begin{exercise}[$\star\star$]\label{exo:g8:powers:8}
Une rumeur se propage~: le jour 1, trois personnes la connaissent~;
chaque jour, chaque personne qui la connaît en informe trois
nouvelles. Écrire avec une puissance le nombre de \emph{nouvelles}
personnes informées le jour $4$, et calculer combien de personnes
connaissent la rumeur à la fin du jour 4 (y compris les trois
d'origine).
\end{exercise}

\begin{exercise}[$\star\star$]\label{exo:g8:powers:9}
Une feuille de papier a $0.1$ mm d'épaisseur, soit $10^{-4}$ m. La
plier double son épaisseur à chaque fois.
\begin{enumerate}
  \item Exprimer l'épaisseur après $10$ plis comme un \omterm{def:g2:mult:def}{produit}, et la
        calculer en centimètres ($2^{10} = 1024$).
  \item Après $42$ plis l'épaisseur serait
        $2^{42} \times 10^{-4}$ m, avec $2^{42} \approx 4.4 \times
        10^{12}$. Montrer que cela dépasse la distance Terre--Lune,
        environ $3.8 \times 10^8$ m.
\end{enumerate}
\end{exercise}

\begin{exercise}[$\star\star$]\label{exo:g8:powers:10}
Ranger du plus petit au plus grand, sans calculatrice~:
\[
2^{10}, \qquad 10^3, \qquad 3^6, \qquad 5^4 .
\]
(Calculer chacun~; $2^{10}$ et $10^3$ sont des voisins célèbres.)
\end{exercise}

\begin{exercise}[$\star\star\star$]\label{exo:g8:powers:11}
Lequel est plus grand, $2^{100}$ ou $10^{30}$~? Utiliser
$2^{10} = 1024 > 10^3$ pour comparer
$2^{100} = \left(2^{10}\right)^{10}$ avec
$\left(10^3\right)^{10}$.
\end{exercise}

\section{Problème~: l'échiquier et les puissances de deux}

\begin{problem}[{Devoir maison --- la \omterm{def:g1:addition:def}{somme} géométrique
$1 + 2 + 4 + \dots + 2^{n-1} = 2^n - 1$, d'une légende célèbre
jusqu'aux nombres binaires}]
\label{pb:g8:powers:1}
La légende~: en récompense d'avoir inventé le jeu d'échecs, le sage
Sissa demanda à son roi un grain de blé sur la première case du
plateau, deux sur la deuxième, quatre sur la troisième --- en
doublant de case en case, jusqu'à la soixante-quatrième. Le roi rit
d'une telle modestie. Ce problème calcule ce que le roi a promis, à
l'aide des règles sur les \omterm{def:g8:powers:def}{exposants} du~\cref{thm:g8:powers:rules}, et
s'achève là où l'histoire mène en secret~: les nombres binaires
qu'abrite tout ordinateur.

\medskip
\noindent\textbf{Partie I --- L'astuce du doublement.}
Pour $n \geq 1$, notons $S_n$ le nombre total de grains posés sur les
$n$ premières cases.
\begin{enumerate}
  \item Exprimer le nombre de grains de la case $k$ comme une
        puissance de $2$. Quelle puissance se trouve sur la case
        $64$~?
  \item Calculer $S_1$, $S_2$, $S_3$, $S_4$ et $S_5$, et comparer
        chacun à une puissance de $2$ voisine. Conjecturer une
        formule pour $S_n$.
  \item L'\emph{astuce du doublement}~: écrire les \omterm{def:g1:addition:def}{sommes} $S_n$ et
        $2 \times S_n$ l'une sous l'autre, soustraire, et démontrer
        la conjecture~:
        \[
        S_n = 1 + 2 + 4 + \dots + 2^{n-1} = 2^n - 1 .
        \]
  \item Combien de grains le roi a-t-il promis en tout~? Exprimer la
        réponse à l'aide d'une puissance de $2$, et compléter la
        remarque classique~: «~l'échiquier entier porte un grain de
        moins qu'une seule soixante-cinquième case n'en porterait~».
  \item Montrer que la seconde \omterm{def:g3:division:half}{moitié} de l'échiquier (cases $33$ à
        $64$) porte \emph{exactement} $2^{32}$ fois plus de grains
        que la première \omterm{def:g3:division:half}{moitié}.
\end{enumerate}

\medskip
\noindent\textbf{Partie II --- Quelle est la taille de $2^{64}$~?}
La comparaison $2^{10} = 1024 > 10^3$ de
l'\cref{exo:g8:powers:11} est la clé de toutes les estimations qui
suivent.
\begin{enumerate}[resume]
  \item Montrer que
        $2^{64} = 2^4 \times \left(2^{10}\right)^6 > 1.6 \times
        10^{19}$.
  \item Un grain de blé pèse environ $0.05$ g, c'est-à-dire
        $5 \times 10^{-2}$ g. Montrer que le blé promis pèse plus de
        $8 \times 10^{17}$ g, et convertir cette \omterm{def:g2:measure:units}{masse} en tonnes
        ($1$ tonne $= 10^6$ g).
  \item Le monde entier récolte actuellement environ
        $8 \times 10^8$ tonnes de blé par an. Combien d'années de
        récolte mondiale le roi a-t-il promis, au moins~?
  \item Trouver le plus petit entier $n$ tel que $2^n > 10^6$ ---
        c'est-à-dire combien de doublements il faut pour dépasser le
        million. (Calculer $2^{19}$ et $2^{20}$ exactement, en
        utilisant $2^{10} = 1024$.)
  \item Un beau parleur nous propose un salaire mensuel~: $1$
        centime le premier jour, puis chaque jour le double de la
        veille. Quel jour la paie \emph{journalière} dépasse-t-elle
        pour la première fois un million d'euros ($10^8$
        centimes)~? (Calculer $2^{26}$ et $2^{27}$ exactement.)
\end{enumerate}

\medskip
\noindent\textbf{Partie III --- Des poids binaires.}
Une marchande possède cinq poids~: $1$, $2$, $4$, $8$ et $16$
grammes, un de chaque. Elle en pose certains sur l'un des plateaux
d'une balance pour peser des marchandises sur l'autre plateau.
\begin{enumerate}[resume]
  \item Quels poids pose-t-elle pour peser $21$ g~? Pour peser
        $27$ g~?
  \item Expliquer pourquoi toute cible de $16$ g ou plus \emph{doit}
        utiliser le poids de $16$ g, et pourquoi toute cible de
        $15$ g ou moins ne doit \emph{pas} l'utiliser. (La question
        3 dit ce que les poids $1, 2, 4, 8$ peuvent atteindre au
        maximum.) Expliquer pourquoi le même raisonnement se répète
        avec le poids suivant par ordre décroissant, à chaque étape.
  \item En déduire que toute cible entière de $1$ à $31$ g peut être
        pesée, et \emph{d'une seule façon}~: chaque nombre entre $1$
        et $31$ est une \omterm{def:g1:addition:def}{somme} de puissances de $2$ distinctes d'une
        unique manière.
  \item La marchande achète un sixième poids, de $32$ g. Jusqu'à
        quelle cible peut-elle désormais peser~? Écrire $45$ g comme
        \omterm{def:g1:addition:def}{somme} de puissances de $2$ distinctes.
  \item Dans un ordinateur, un nombre «~sur 64 bits~» est stocké sur
        $64$ cases contenant chacune un $0$ ou un $1$ --- la case
        $k$ apportant $2^{k-1}$ lorsqu'elle contient un $1$, comme
        les grains de la légende. À l'aide de la partie I, expliquer
        pourquoi les nombres entiers qu'une telle machine peut
        stocker vont exactement de $0$ à $2^{64} - 1$.
\end{enumerate}
\end{problem}
